\documentclass[10pt,a4paper]{amsart}
\usepackage[margin=19mm]{geometry}
\usepackage{amsmath,amssymb,mathtools}
\usepackage[scr=rsfs,cal=euler]{mathalfa}
\usepackage{xfrac}
\usepackage[unicode,hidelinks,bookmarks=false]{hyperref}
\newcommand{\ie}{i.e.\ }
\newcommand{\cf}{cf.\ }
\newcommand{\resp}{respectively\ }
\newcommand{\Subsection}{\subsection}
\providecommand{\nobreakdash}{\nobreak\hbox{-}}
\setlength{\emergencystretch}{2em}
\multlinegap0pt
\usepackage{bm}
\usepackage{bbm}
\usepackage{dsfont}
\usepackage{xspace}
\usepackage{cancel}
\usepackage{mathtools}
\usepackage{amsthm}
\makeatletter
\@ifundefined{theorem}{\newtheorem{theorem}{Theorem}}{}
\makeatother
\makeatletter
\expandafter\ifx\csname audit\endcsname\relax
\newenvironment{audit}
  {\par\smallskip\noindent\textbf{Audit points.}\begin{itemize}}
\fi
\newcommand{\nat}{\mathbb N}
\makeatother
\title[The Algebra of compact-open subsets in the spectrum of the r]{The Algebra of compact-open subsets in the spectrum of the ring $C(T)$ for an infinite compact Hausdorff space $T$}
\author{Evgeny Kuznetsov}
\date{Source: arXiv:2608.08251v1 (CC BY 4.0)}
\begin{document}
\maketitle

\noindent\emph{Source and licence.} The Algebra of compact-open subsets in the spectrum of the ring $C(T)$ for an infinite compact Hausdorff space $T$. Version arXiv:2608.08251v1 (CC BY 4.0), distributed under the Creative Commons Attribution 4.0 International licence, as recorded in the arXiv licence metadata for this exact version and shown on the abstract page https://arxiv.org/abs/2608.08251v1. Full rights notice: licenses/arxiv-2608.08251-CC-BY-4.0.txt.\par\smallskip

\noindent\textit{Editorial scope.} Counted unit: the Theorem whose source label is \texttt{thm:grid-existence} in the article (source0.tex lines 1091--1097, source label \texttt{thm:grid-existence}), delivered together with the article's own complete original proof of it (source0.tex lines 1099--1201, one \texttt{proof} environment of 103 lines). No mathematics is written, repaired or re-derived: statement and proof are the article's own text line for line. Required context retained uncounted: none. Every definition, display and label of those lines is retained unchanged. Omitted from this package and not redistributed: 1--1090, 1098--1098, 1202--1338. Nothing inside a retained excerpt is omitted or altered: every retained body is delivered exactly as the article's lines are written, so no omission receipt of any kind accompanies this package. Citations: the retained excerpts carry no citation command at all, so no bibliography entry of the article is transcribed and no reference is redistributed. Reference adaptation: none was needed and none was made; every reference inside the delivered unit resolves inside it, so no unresolved reference or citation is produced. Printed slips kept literal: no printed word, symbol, hypothesis or number of the counted statement or of its proof was altered. Layout and class: the article's own class is \texttt{amsart} with packages including fontenc, lmodern, amsmath, amssymb, amsthm, mathtools, biblatex; the wrapper is the lane's pinned \texttt{amsart} editorial wrapper at 10pt on A4, which restates the theorem-like environments the retained text needs and adds the article's own macro definitions that the retained text needs (\texttt{\textbackslash nat}), with extra wrapper packages (bm, bbm, dsfont, xspace, cancel, mathtools). The delivered numbering is the unit's own. Assets: no retained span contains a figure environment, a graphics-inclusion command, a PDF-page inclusion, a table or a TikZ picture; no image file is copied into this package, the package declares no asset dependency and no image file is redistributed with it. Licence: version arXiv:2608.08251v1 of this article is distributed under the Creative Commons Attribution 4.0 International licence, as recorded in the arXiv licence metadata for this exact version and shown on the abstract page https://arxiv.org/abs/2608.08251v1. A full rights notice is delivered at licenses/arxiv-2608.08251-CC-BY-4.0.txt. Characters: every retained excerpt is pure ASCII, so the delivered engine, XeLaTeX with its default Latin Modern text font, reports no missing character.
\begin{theorem}\label{thm:grid-existence}
Every infinite $\sigma$-complete Boolean algebra $\mathcal B$ admits a
countable grid: a partition of unity $\{B_j\}_{j\geq1}$ (pairwise
disjoint, nonzero, $\bigvee_jB_j=1$) together with, for each $j$, a
countable partition $\{C_{j,t}\}_{t\geq1}$ of $B_j$ (pairwise disjoint,
nonzero, $\bigvee_tC_{j,t}=B_j$).
\end{theorem}

\begin{proof}
\emph{Step 0: an unconditional antichain.} Since $\mathcal B$ is infinite
it has more than two elements; pick $0<c_1<1$ and set $d_1:=c_1^c$ (both
nonzero). Since $\mathcal B\cong{\downarrow}c_1\times{\downarrow}d_1$ and
$\mathcal B$ is infinite, at least one factor is infinite; relabel so
that ${\downarrow}c_1$ is infinite. Recursing inside ${\downarrow}c_1$:
pick $0<c_2<c_1$ with ${\downarrow}c_2$ infinite, and set
$d_2:=c_1\wedge c_2^c$. Iterating gives a strictly decreasing sequence
$c_1>c_2>\cdots$ with ${\downarrow}c_n$ infinite at each stage, and
$d_n:=c_{n-1}\wedge c_n^c$ (with $c_0:=1$) pairwise disjoint (if $m>n$
then $d_m\leq c_{m-1}\leq c_n$, so
$d_m\wedge d_n\leq c_n\wedge c_{n-1}\wedge c_n^c=0$) and nonzero by
construction. This gives an antichain $\{d_n\}_{n\geq1}$, valid in any
infinite Boolean algebra, with no completeness assumption.

\emph{Step 1: countable distributivity.} For $x\in\mathcal B$ and
countable $\{y_n\}\subseteq\mathcal B$,
\[
x\wedge\bigvee_ny_n=\bigvee_n(x\wedge y_n)
\]
(both joins exist by $\sigma$-completeness). Writing $s=\bigvee y_n$,
$s'=\bigvee(x\wedge y_n)$: $s'\leq x\wedge s$ is immediate. For the
converse, $x\wedge s\leq s'$ is equivalent to $s\leq s'\vee x^c$; and for
each $n$, $y_n=(y_n\wedge x)\vee(y_n\wedge x^c)\leq s'\vee x^c$ (since
$y_n\wedge x\leq x\wedge y_n\leq s'$), so $s'\vee x^c$ is an upper bound
of $\{y_n\}$, giving $s\leq s'\vee x^c$ since $s$ is the least such.
Hence $x\wedge s\leq s'$, so $x\wedge s=s'$.

\emph{Step 2: two cases by number of atoms.} Let $\mathrm{At}(\mathcal
B)$ be the set of atoms of $\mathcal B$.

\emph{Case A: $\mathrm{At}(\mathcal B)$ infinite.} Choose a countably
infinite subset $\{p_1,p_2,\dots\}\subseteq\mathrm{At}(\mathcal B)$ and
partition it into countably many countably infinite pieces
$P_1,P_2,\dots$ via any bijection $\nat\cong\nat\times\nat$. By
$\sigma$-completeness, $B_j:=\bigvee_{p\in P_j}p$ exists. Then $B_j\neq0$
(it dominates any $p\in P_j$); for $j\neq k$,
\[
B_j\wedge B_k=\bigvee_{p\in P_j}(p\wedge B_k)
=\bigvee_{p\in P_j}\bigvee_{q\in P_k}(p\wedge q)=0
\]
by Step~1 (distinct atoms are disjoint); ${\downarrow}B_j$ is infinite,
since distinct finite $F\neq F'\subseteq P_j$ give distinct joins
$\bigvee F\neq\bigvee F'$ (an atom $p\in P_j$ satisfies $p\leq\bigvee F$
iff $p\in F$, again by Step~1), and $P_j$ is countably infinite; and the
atoms of ${\downarrow}B_j$ are exactly $P_j$ (if $q$ is an atom
$\leq B_j$, then $q=q\wedge B_j=\bigvee_{p\in P_j}(q\wedge p)\neq0$, so
$q\wedge p\neq0$ for some $p\in P_j$, forcing $q=p$). Finally,
${\downarrow}B_j$ is $\sigma$-complete: any countable
$\{y_n\}\subseteq{\downarrow}B_j$ has $\bigvee y_n$ (computed in
$\mathcal B$) satisfying $\bigvee y_n\leq B_j$, since $B_j$ is an upper
bound of $\{y_n\}$ and $\bigvee y_n$ is the least such. So
${\downarrow}B_j$ is again an infinite $\sigma$-complete Boolean algebra
with countably many atoms; Case~A recurses inside it to produce a
countably infinite antichain below $B_j$, each member again of infinite
relativization.

\emph{Case B: $\mathrm{At}(\mathcal B)$ finite (possibly empty).} Let
$e:=\bigvee\mathrm{At}(\mathcal B)$, a finite join. A direct computation
($x=\bigvee_i(x\wedge p_i)$ for $x\leq e$, the $p_i$ the finitely many
atoms) shows ${\downarrow}e$ consists exactly of the sub-joins of the
$p_i$, so it is finite. Since $\mathcal B\cong{\downarrow}e\times
{\downarrow}e^c$ and $\mathcal B$ is infinite, ${\downarrow}e^c$ is
infinite; and since every atom of $\mathcal B$ is $\leq e$,
${\downarrow}e^c$ is atomless. In an atomless algebra every nonzero
$x\leq e^c$ is itself atomless (an atom $\leq x$ would be an atom
$\leq e^c$) and hence, being a nonzero atomless element, never an atom
and always of infinite relativization (a finite nonzero Boolean algebra
is always atomic). So Step~0's peeling, run inside ${\downarrow}e^c$,
produces a countably infinite antichain $\{d_n\}$ below $e^c$, every
member of which is again atomless (hence again infinite), and the same
argument recurses with no case distinction needed. This sub-case uses no
$\sigma$-completeness, only finitely many splits repeated countably
often.

\emph{Step 3: from antichain to partition.} Steps~0--2 produce, below
any given $e\in\mathcal B$ (with $e=1$ at the top level, or $e=B_j$ when
recursing inside ${\downarrow}B_j$), a countably infinite antichain
$\{E_n\}_{n\geq1}$ of pairwise disjoint nonzero elements each of
infinite relativization --- but the antichain need not itself join to
$e$. To repair this, set
\[
s:=\bigvee_nE_n\quad(\text{exists by }\sigma\text{-completeness}),
\qquad
r:=e\wedge s^c,
\]
and define $E_1':=E_1\vee r$, $E_n':=E_n$ for $n\geq2$. Since $r\leq
s^c$ is disjoint from every $E_n\leq s$, and the $E_n$ are themselves
pairwise disjoint, $\{E_n'\}$ is again pairwise disjoint; each $E_n'$ is
nonzero and has infinite relativization (as $E_n'\geq E_n$). Since
$E_n\leq e$ for every $n$, also $s\leq e$, so $e\wedge s=s$ and
\[
\bigvee_nE_n'=r\vee\bigvee_nE_n=r\vee s=(e\wedge s^c)\vee(e\wedge s)=e.
\]
So $\{E_n'\}$ is a genuine countable partition of $e$ into pieces of
infinite relativization.

Applying Step~3 with $e=1$ to the antichain from Step~2 gives the
partition of unity $\{B_j\}_{j\geq1}$. Applying Step~3 again, with
$e=B_j$, to the antichain that Step~2 recurses to inside
${\downarrow}B_j$, gives the partition $\{C_{j,t}\}_{t\geq1}$ of $B_j$.
This is the required grid.
\end{proof}

\end{document}
